%%% preamble.tex
%
% This is where one loads LaTeX packages and defines their own commands. Some
% of the commands you see here are defined in the file
% presentation-metadata.tex.
%

%% Load packages.

\usepackage[utf8]{inputenc}

\usepackage[\mylanguage]{babel}

\usepackage{fourier}

\usepackage{mathtools}

\usepackage[\citationsettings]{biblatex}

\addbibresource{../content/sources.bib}

\usepackage{csquotes}

%% Commands.

\NewDocumentCommand\bvec{ m }{\vec{\mathbf{#1}}}

\DeclareMathOperator\divop{div}

\DeclarePairedDelimiter\arcs{(}{)}

%% Set graphicspath to contents/images/, so users don't need to enter the full path of the file.

\graphicspath{{../content/images/}}

%% Patch the item separation between itemize and enumerate items.

\usepackage{xpatch}

\xpatchcmd
  {\itemize}
  {\def\makelabel}
  {\setlength{\itemsep}{1ex}\def\makelabel}
  {}
  {}

\xpatchcmd
  {\enumerate}
  {\def\makelabel}
  {\setlength{\itemsep}{1ex}\def\makelabel}
  {}
  {}

%% Define inline verbatim command with xparse.

\NewDocumentCommand\code{m}{\texttt{#1}}

\NewDocumentCommand\latexcommand{m}{\textbackslash\code{#1}}

%% Define temporary image background style command.
%
% Invoke before a frame to set a background image for the following frames. If
% a background image is to be set for a single frame, surround this command
% and the frame itself with \begingroup … \endgroup.
%
\NewDocumentCommand\setbackgroundimage{
    O{0.9}   % Optional image width in relation to paper width, given inside [⋅].
    D(){0.5} % Optional image x position in relation to paper height, given inside (⋅).
    D<>{0.4} % Optional image y position in relation to paper height, given inside <⋅>.
    m        % A mandatory path to the background image.
}{
    \setbeamertemplate{background}{
        \begin{tikzpicture}
            \useasboundingbox (0,0) rectangle (\the\paperwidth,\the\paperheight);
            \fill[color=white] (0,0) rectangle (\the\paperwidth,\the\paperheight);
            \node at (\taulogox, \taulogoy) {\includegraphics[width=\taulogow]{../tau-logo/tau-logo.pdf}};
            \node at (#2\paperwidth, #3\paperheight) {\includegraphics[width=#1\paperwidth]{#4}};
        \end{tikzpicture}
    }
}

%% A command for a final unnumbered page.

\NewDocumentCommand\finalpage{m}{
    \begingroup
        \setbeamertemplate{background}{
            \begin{tikzpicture}
                \useasboundingbox (0,0) rectangle (\the\paperwidth,\the\paperheight);
                \fill[color=TAUPURPLE] (0,0) rectangle (\the\paperwidth,\the\paperheight);
            \end{tikzpicture}
        }
        \begin{frame}[noframenumbering]{}
            \vfill
            \centering
            \color{TAUWHITE}
            \resizebox{0.6\textwidth}{!}{#1}
        \end{frame}
    \endgroup
}

%% Citation bracket color.

\DeclareOuterCiteDelims{cite}{\textcolor{TAUPURPLE}{\bibopenbracket}}{\textcolor{TAUPURPLE}{\bibclosebracket}}

%% Write metadata to an external file, where pdfx can read it from.

\begin{filecontents*}[overwrite]{\jobname.xmpdata}
\Title{\mytitle}
\Author{\myauthor}
\Subject{\mysubject}
\Keywords{\mykeywords}
\Creator{\mycreator}
\end{filecontents*}

%% PDF/A

\usepackage[a-3a, mathxmp]{pdfx}

\hypersetup{
    colorlinks,
    allcolors=black,
    urlcolor=TAUPURPLE,
    citecolor=TAUPURPLE,
    anchorcolor=TAUPURPLE
}

\usepackage{url}
